Bij langere packet verzendtijd: Als het slot langer is dan de pakketduur, wordt het pakket succesvol verzonden, maar blijft het kanaal aan het einde van elk slot een tijdje onnodig stil (dead time). Omdat genormaliseerde throughput de fractie van de totale tijd is waarin succesvolle data wordt overgedragen, zorgt deze loze ruimte voor een efficiëntieverlies. De curve behoudt dezelfde vorm, maar de maximale throughput zal lineair dalen naarmate het slot groter wordt, omdat er simpelweg capaciteit wordt verspild aan stilte.\\

Bij kortere packet verzendtijd: Als een slot korter is dan de tijd die nodig is om het pakket volledig te verzenden, "bloedt" een transmissie die start in slot k door in slot k+1. Als een ander station in slot k+1 besluit te gaan zenden, ontstaat er alsnog een botsing halverwege het pakket. De synchronisatie is hiermee kapot, en het netwerkgedrag degenereert direct terug naar dat van pure Aloha